% Generated draft: requires mathematical and source review.
\documentclass[11pt]{article}
\usepackage[margin=25mm]{geometry}
\usepackage{fontspec}
\setmainfont{Latin Modern Roman}
\usepackage{amsmath,amssymb,mathtools,unicode-math}
\setmathfont{Latin Modern Math}
\usepackage{enumitem,needspace,xurl,hyperref}
\setlength{\parindent}{0pt}
\setlength{\parskip}{5pt}
\setlength{\emergencystretch}{3em}
\newcommand{\N}{\mathbb{N}}
\newcommand{\Z}{\mathbb{Z}}
\newcommand{\Q}{\mathbb{Q}}
\newcommand{\R}{\mathbb{R}}
\newcommand{\C}{\mathbb{C}}
\newcommand{\F}{\mathbb{F}}
\newcommand{\A}{\mathbb{A}}
\newcommand{\PP}{\mathbb{P}}
\newcommand{\E}{\mathbb{E}}
\newcommand{\eps}{\varepsilon}
\newcommand{\dd}{\,\mathrm{d}}
\begin{document}
\section*{UnsolvedMath ID: 30001239}
\textbf{Source problem code:} catalogue entry
\subsection{Definitions and mathematical statement}
Let \(k\) be a field and let \(S=k[x_1,\ldots,x_n]\) have its standard grading, with \(\deg x_j=1\). Let \(I\subseteq S\) be a homogeneous ideal, put \(R=S/I\), and write \(d=\dim R\). Let \(r\geq1\) be an integer. No restriction on the characteristic of \(k\) is supplied. The graded canonical module of \(S\) is \(\omega_S=S(-n)\), where \(S(-n)_j=S_{j-n}\). For \(0\leq i\leq n\), put
\[
K^i(R)=\operatorname{Ext}^{n-i}_S(R,\omega_S),
\]
using graded Ext with its natural grading. For a nonzero finitely generated graded \(S\)-module \(M\), let \(\dim M\) denote its Krull dimension and define its Castelnuovo--Mumford regularity by
\[
\operatorname{reg}M=\max\{j-p:\operatorname{Tor}^S_p(M,k)_j\ne0\}.
\]
The ring \(R\) satisfies Serre's condition \((S_r)\) if
\[
\operatorname{depth}R_{\mathfrak p}\geq\min\{r,\dim R_{\mathfrak p}\}
\qquad\text{for every }\mathfrak p\in\operatorname{Spec}R.
\]
A monomial ideal is an ideal generated by monomials in \(x_1,\ldots,x_n\).

\emph{Source question, with unresolved qualifications retained.} Determine which homogeneous ideals satisfy the regularity bound
\[
\operatorname{reg}K^i(R)\leq\dim K^i(R).
\tag{A}
\]
The supplied record does not specify the range of indices or the quantifier on \(i\) in this first question. It further asks whether, when \(R\) satisfies \((S_r)\), the corresponding stronger bounds hold:
\[
\operatorname{reg}K^i(R)\leq i-r
\qquad\text{for every integer }0\leq i\leq d-1.
\tag{B}
\]
The question explicitly includes the case of monomial ideals. If \(d=0\), the displayed index range in (B) is empty.

The antecedent of ``they'' in the source is unresolved: (B) may concern the ideals satisfying the intended version of (A), or a broader class of homogeneous ideals whose quotients satisfy \((S_r)\). The supplied background describes an extension beyond classes for which a regularity theorem was already proved, but does not identify those classes or state that theorem. Thus (A) and (B) retain the source's question form and do not assert universal inequalities. The formulas above are unambiguous for nonzero Ext modules. The source does not specify conventions for zero modules; declaring their regularity and dimension to be \(-\infty\) would make the corresponding inequalities automatic, but this is an additional convention requiring review. The source also does not explicitly exclude \(I=S\), for which the usual finite-dimensional quotient formulation needs qualification.
\subsection{Short English statement}
The problem asks which homogeneous ideals have Ext modules whose regularity is at most their dimension, and whether Serre's condition \((S_r)\) yields the bound \(i-r\) for every \(0\leq i\leq d-1\), including for monomial ideals. The first index range, the intended class in the second question, and zero-module conventions remain unresolved. The supplied record labels the problem open; this draft does not independently verify that status.
\subsection{Sources}
\begin{itemize}
\item Local UnsolvedMath registry, supplied complete record with permanent ID \(30001239\), exact source problem code \texttt{catalogue entry}, title ``Regularity Bounds for Ext Modules Under Serre Conditions,'' in the Oberwolfach Reports Open Problems collection. The record gives proposed year 2009 and names no proposer.
\item Supplied report provenance: \emph{Kommutative Algebra} (2009), Oberwolfach Reports, DOI 10.4171/owr/2009/22, pp. 1157--1224. \url{https://doi.org/10.4171/owr/2009/22}. The supplied background describes these reports as documenting mathematical workshops and their open-problem discussions. No location within this page range is supplied for the particular question.
\item Supplied literature reference: ``Murai-Terai on H-vectors and Serre conditions.'' \url{https://arxiv.org/abs/0912.1089}. The record does not identify a particular result in this reference as resolving or precisely delimiting the question.
\item Supplied literature assessment, checked on exactly \texttt{2026-08-21}: status ``open.'' The assessment describes the extracted sentence as a fragment asking for an extension beyond already proved classes, rather than a resolved theorem, and states that current literature gives no general answer. This is supplied dated curation derived from the OWR record, not an independently verified literature assessment; its status and cited claims require independent verification.
\item The registry supplies statement status \texttt{corrected\_verified} and reports comparison with the official Oberwolfach report, with notation, context, and question form normalized without changing the source-level claim. This verification is reported provenance, not a source check performed for this draft.
\end{itemize}
\end{document}
